\begin{derivation}[从\eqnrefs{eq:qnm_left_right_null_vectors}到\eqnrefs{eq:qnm_green_residue_operator_formula}]
设开放 Maxwell 算符束
\[
\mathsf L_{\rm EM}(\omega)
\]
在复频率 \(\widetilde\omega_c\) 处只有一个简单零本征值，并取相应右、左零矢量
\[
\mathsf L_c|f_c^R\rangle=0,
\qquad
\langle f_c^L|\mathsf L_c=0,
\qquad
\mathsf L_c\equiv\mathsf L_{\rm EM}(\widetilde\omega_c).
\]
“简单”意味着左右零空间都为一维，而且
\(\langle f_c^L|\mathsf L_c'|f_c^R\rangle\neq0\)，其中
\(\mathsf L_c'\equiv\partial_\omega\mathsf L_{\rm EM}(\widetilde\omega_c)\)。

令
\[
\delta\omega\equiv\omega-\widetilde\omega_c.
\]
算符束与其逆算符分别作 Taylor--Laurent 展开：
\begin{align*}
\mathsf L_{\rm EM}(\omega)
&=\mathsf L_c+\delta\omega\,\mathsf L_c'
+O(\delta\omega^2),\\
\mathbf G^R(\omega)
&=\frac{\mathsf R_{-1}}{\delta\omega}
+\mathsf R_0+O(\delta\omega).
\end{align*}
这里 \(\mathsf R_{-1}\) 是要确定的留数算符，并未预先假定其形式。

先把两式代入左逆关系
\[
\mathsf L_{\rm EM}(\omega)\mathbf G^R(\omega)=\mathbf I.
\]
\(\delta\omega^{-1}\) 阶给
\[
\mathsf L_c\mathsf R_{-1}=0.
\]
因此 \(\mathsf R_{-1}\) 的每一个输出向量都属于
\(\ker\mathsf L_c=\operatorname{span}\{|f_c^R\rangle\}\)，故必存在某个线性泛函
\(\langle\alpha|\) 使
\[
\mathsf R_{-1}=|f_c^R\rangle\langle\alpha|.
\]
再把同一 Laurent 展开代入右逆关系
\[
\mathbf G^R(\omega)\mathsf L_{\rm EM}(\omega)=\mathbf I.
\]
其 \(\delta\omega^{-1}\) 阶给
\[
\mathsf R_{-1}\mathsf L_c=0.
\]
代入上式后有
\[
|f_c^R\rangle\langle\alpha|\mathsf L_c=0.
\]
因为 \(|f_c^R\rangle\neq0\)，所以
\(\langle\alpha|\mathsf L_c=0\)。左零空间也是一维，因此
\[
\langle\alpha|=\frac{1}{\mathcal N_c}\langle f_c^L|
\]
对某个尚待确定的标量 \(\mathcal N_c\)。由此，秩一结构不是猜出来的，而是同时由左右逆关系强制得到：
\[
\mathsf R_{-1}
=\frac{|f_c^R\rangle\langle f_c^L|}{\mathcal N_c}.
\]

现在比较 \(\mathsf L\mathbf G=\mathbf I\) 的常数阶：
\[
\mathsf L_c\mathsf R_0
+\mathsf L_c'\mathsf R_{-1}
=\mathbf I.
\]
左乘 \(\langle f_c^L|\)。第一项因为
\(\langle f_c^L|\mathsf L_c=0\) 消失，于是
\begin{align*}
\langle f_c^L|
&=\langle f_c^L|\mathsf L_c'\mathsf R_{-1}\\
&=\frac{\langle f_c^L|\mathsf L_c'|f_c^R\rangle}
{\mathcal N_c}\langle f_c^L|.
\end{align*}
由于左零矢量非零，比较其系数得到
\[
\boxed{
\mathcal N_c
=\langle f_c^L|
\partial_\omega\mathsf L_{\rm EM}(\widetilde\omega_c)
|f_c^R\rangle }.
\]
因此 Green 算符在该简单共振附近的唯一极点部分为
\[
\boxed{
\mathbf G^R(\omega)
=\frac{|f_c^R\rangle\langle f_c^L|}
{(\omega-\widetilde\omega_c)
\langle f_c^L|\partial_\omega\mathsf L_{\rm EM}(\widetilde\omega_c)|f_c^R\rangle}
+\mathbf G_{\rm reg}^R(\omega)}.
\]
这就是正文\eqnrefs{eq:qnm_green_residue_operator_formula}。若
\(|f_c^R\rangle\mapsto a|f_c^R\rangle\)、
\(\langle f_c^L|\mapsto a^{-1}\langle f_c^L|\)，分子与分母同时不变，所以 Green 留数不依赖左右本征矢量各自的任意归一化。

闭合、无损、非色散情形可作为检查。此时
\[
\partial_\omega\mathsf L_{\rm EM}
=-\frac{2\omega}{c^2}\boldsymbol\varepsilon_r.
\]
利用正文\eqnrefs{eq:lossless_electric_mode_normalization}，有
\[
\mathcal N_c=-\frac{2\omega_c}{c^2},
\qquad
\mathsf R_{-1}
=-\frac{c^2}{2\omega_c}
\mathbf f_c\otimes\mathbf f_c^*.
\]
另一方面，离散模谱表示中的频率因子可以直接部分分式分解：
\begin{align*}
\frac{c^2\mathbf f_c\otimes\mathbf f_c^*}{\omega_c^2-\omega^2}
&=-\frac{c^2}{2\omega_c}
\frac{\mathbf f_c\otimes\mathbf f_c^*}{\omega-\omega_c}
+\frac{c^2}{2\omega_c}
\frac{\mathbf f_c\otimes\mathbf f_c^*}{\omega+\omega_c}.
\end{align*}
第二项在 \(\omega=\omega_c\) 解析，因此第一项的系数正是上面由算符束推得的留数。开放 QNM 归一化由此与熟悉的闭合模归一化连续接上。
\end{derivation}
